\section{A causal stationary field on the original exterior}
\label{q4:adj:section}

We continue with the original equality data $(\Omega,g,K)$ of
Definition~\ref{q4:def:horizon}.  In particular, $S=\partial\Omega$ is connected,
$\theta_+=0$ on $S$, and $S$ is outer area-minimizing.  The constants fixed in
Section~\ref{q4:var:section} are
\[
 r_0=\sqrt{2m},\qquad b^0=\frac E m,\qquad b^i=-\frac{P_i}{m},
 \qquad c=\frac2{r_0},\qquad \kappa=\frac c2=\frac1{r_0}.
\]
Thus $(b^0)^2-|b|^2=1$ and $b^0>0$.  We fix the same exponent
$1<q_0<\min\{q,2\}$ as in the strict variation constructed there.
The purpose of this section is to extract a smooth field from the multiplier
and to retain all information about the original second fundamental form.

\begin{theorem}[Causal stationary field]
\label{q4:adj:causal-field}
There are a smooth function $u$ and a smooth vector field $X$ on $\Omega$,
smooth up to its one-sided boundary, with the following properties.
\begin{enumerate}[label=\textup{(\roman*)}]
\item They satisfy $u\ge |X|_g$, $u>0$ in $\operatorname{int}\Omega$, and
\begin{equation}
 u\mu+J(X)=0.
 \label{q4:adj:complementarity}
\end{equation}
With symmetric indices denoting averaging, their equations are
\begin{align}
 \nabla_{(i}X_{j)}&=-uK_{ij},
 \label{q4:adj:kid-first}\\
 \Delta u&=-\operatorname{div}_g\bigl(K(X,\cdot)^\sharp\bigr),
 \label{q4:adj:conformal-adjoint}\\
 \nabla^2u
 &=u(\Ric_g+\tau K-2K^2)-\mathcal L_XK+X_{(i}J_{j)}.
 \label{q4:adj:kid-second}
\end{align}
Here $X_i=g_{ij}X^j$ and $(K^2)_{ij}=K_{ik}K^k{}_j$.
\item On the chosen asymptotically flat end,
\begin{equation}
 (u,X)=(b^0,b)+O_2(r^{-q_0}).
 \label{q4:adj:decay}
\end{equation}
\item At the boundary, for the unit normal $\nu$ pointing into $\Omega$,
\begin{equation}
 X=u\nu,\qquad
 \partial_\nu u+K(X,\nu)=\kappa.
 \label{q4:adj:horizon-data}
\end{equation}
Either $u>0$ everywhere on $S$ or $u=0$ everywhere on $S$.
\item On $\mathbb R\times\operatorname{int}\Omega$, set
\begin{equation}
 \mathbf g=-u^2\,\mathrm dt^2
   +g_{ij}(\mathrm dx^i+X^i\,\mathrm dt)
          (\mathrm dx^j+X^j\,\mathrm dt),
 \qquad \xi=\partial_t,\qquad N=u^2-|X|_g^2.
 \label{q4:adj:stationary-metric}
\end{equation}
For the future normal $n=u^{-1}(\xi-X)$, the slice $t=0$ has exactly the
original induced metric $g$ and second fundamental form $K$, with the convention
of Theorem~\ref{q4:thm:rigidity}.  Moreover,
\begin{equation}
 \Ric_{\mathbf g}(\xi,\cdot)=0,
 \qquad \Ric_{\mathbf g}=0\quad\hbox{on }\{N>0\}.
 \label{q4:adj:ricci}
\end{equation}
\item We have $N\ge0$, $N\to1$ at infinity, and $N>0$ in a full collar just
outside $S$.  More precisely, in the positive boundary-lapse case,
\begin{equation}
 N|_S=0,\qquad
 \mathrm dN|_S=2u\kappa\,\nu^\flat=2\kappa X^\flat.
 \label{q4:adj:positive-collar}
\end{equation}
In the zero boundary-lapse case, if $s\ge0$ is the original $g$-normal distance
from $S$, there are smooth one-sided fields $d,Y_2$ such that
\begin{equation}
 u=sd,\qquad X=s^2Y_2,\qquad d|_S=\kappa,\qquad
 N=s^2\bigl(d^2-s^2|Y_2|_g^2\bigr).
 \label{q4:adj:zero-collar}
\end{equation}
Consequently all zeros of $N$ in $\operatorname{int}\Omega$ lie in a compact
subset of that interior.
\end{enumerate}
\end{theorem}

The proof occupies the remainder of this section.  At no point do we assume
that the original data are vacuum or that they lie in a prescribed stationary
development.  The metric in \eqref{q4:adj:stationary-metric} is constructed from
the multiplier.

\subsection{From a positive measure to the interior adjoint equations}
\label{q4:adj:measures}

By Proposition~\ref{q4:var:multiplier} and the area localization in
Proposition~\ref{q4:var:area-localization}, the multiplier identity is
\begin{equation}
 6\omega\,\mathcal E'(H)-c\,a'_S(h)
 =\langle C'_H,\Lambda\rangle
   -\int_S\theta'_H\,\mathrm d\eta+za,
 \qquad H=aQ+H_0.
 \label{q4:adj:localized-multiplier}
\end{equation}
The coefficient $a$ is the coefficient of the strict direction $Q$ in the
variation space.  We recall that
\[
 C(x,v)=R_g+\tau^2-|K|^2+2J(v),\qquad |v|_g\le1,
\]
that $\Lambda$ is a positive locally finite measure on the active set $C=0$,
and that the test-vector identification for a metric variation $h$ has
$\dot v=-\tfrac12h^\sharp v$.

Let $u$ be the scalar pushforward of $\Lambda$ to $\Omega$, and let $X$ be its
vector first moment.  Initially these symbols denote measures: for a compactly
supported scalar $f$ and covector $\alpha$,
\[
 \langle u,f\rangle=\int f(x)\,\mathrm d\Lambda(x,v),
 \qquad
 \langle X,\alpha\rangle=\int\alpha_x(v)\,\mathrm d\Lambda(x,v).
\]
The unit-ball condition and positivity give $|X|_g\le u$ as measures, while
activity gives $\mu u+J(X)=0$ as a measure.  We use $\mathrm dV_g$ as the
reference density when expressing the resulting distributional equations.
After establishing smoothness, we use the same letters for the smooth
densities of these measures.

\begin{lemma}[Interior adjoint equations]
\label{q4:adj:interior-equations}
The moment measures have smooth densities on $\operatorname{int}\Omega$.
They satisfy \eqref{q4:adj:complementarity}--\eqref{q4:adj:kid-second} there, and
$u\ge|X|_g$ there.
\end{lemma}

\begin{proof}
For a variation supported in the open exterior, every term of
\eqref{q4:adj:localized-multiplier} except the constraint pairing vanishes.
For a pure tensor variation $p=\dot K$, that pairing is
\[
 2\bigl\langle u,\tau\tr_g p-K:p\bigr\rangle
 +2\bigl\langle X,\nabla^j(p_{ij}-(\tr_g p)g_{ij})\bigr\rangle.
\]
Its distributional adjoint equation is
\begin{equation}
 u(\tau g-K)-\nabla_{(i}X_{j)}+(\operatorname{div}X)g=0.
 \label{q4:adj:tensor-adjoint}
\end{equation}
Taking the trace in four dimensions gives
$3u\tau+3\operatorname{div}X=0$.  Substitution proves
\eqref{q4:adj:kid-first}.

For the conformal variation $(h,p)=(2\psi g,\psi K)$, the constraint
transformation already established in Section~\ref{q4:var:section} gives,
on active rays,
\[
 C'_H(x,v)=6\{-\Delta\psi+K(\nabla\psi,v)\}.
\]
Its adjoint is \eqref{q4:adj:conformal-adjoint}.  This calculation is valid for
measure moments, so both equations hold in distributions before any
regularity assertion.

Diverge \eqref{q4:adj:kid-first} and substitute
$\operatorname{div}X=-u\tau$.  Commuting covariant derivatives gives a Laplace
equation for $X$ whose right side consists of smooth coefficients times
$u,X,\nabla u$.  In \eqref{q4:adj:conformal-adjoint}, the right side consists of
smooth coefficients times $X,\nabla X$.  Thus the coupled system has diagonal
Laplace principal part on the scalar and vector bundles; all couplings have
order at most one.  Distributional interior elliptic regularity, iterated
with the smooth coefficients, makes both densities smooth.  The measure
inequalities and complementarity now give the asserted pointwise statements.

It remains to verify the full metric equation, including its matter term.
We may now work with smooth $u,X$.  For a pure metric variation $h$, keep
covariant $K$, scalar $u$, and contravariant $X$ fixed.  The constraint pairing
is the variation of
\begin{equation}
 \int\bigl[u(R+\tau^2-|K|^2)
      +2X^i\nabla^j(K_{ij}-\tau g_{ij})\bigr]\,\mathrm dV_g
 \label{q4:adj:metric-functional}
\end{equation}
minus $\int h(X,J^\sharp)\,\mathrm dV_g$.
The latter is precisely the contribution of $\dot v=-\tfrac12h^\sharp v$.
In using the varying volume form in \eqref{q4:adj:metric-functional}, no term
has been added: its original integrand is
$2u\mu+2J(X)=0$ pointwise.

For completeness, put $\pi^{ij}=K^{ij}-\tau g^{ij}$.  Integrating the shift
term once by parts gives
$-\int\pi^{ij}(\mathcal L_Xg)_{ij}\,\mathrm dV_g$, with unchanged boundary
terms outside the variation support.  At fixed covariant $K$,
\[
 \delta\pi^{ij}
 =-h^i{}_aK^{aj}-h^j{}_aK^{ia}
   +(K:h)g^{ij}+\tau h^{ij}.
\]
Variation of the last integral, followed by integration of the term
$-\pi^{ij}(\mathcal L_Xh)_{ij}$, has coefficient
\begin{equation}
 \mathcal L_XK-(\operatorname{div}X)K
  -\{X(\tau)+K^{ab}\nabla_aX_b\}g.
 \label{q4:adj:shift-coefficient}
\end{equation}
One can see the cancellations directly: if $B=\mathcal L_Xg$, the variation
of $\pi$ contributes
$2\Sym(KB)-(\tr B)K-\tau B$; integration of $\mathcal L_Xh$ contributes
$(\mathcal L_X\pi)^{ij}+(\operatorname{div}X)\pi^{ij}$ with indices lowered;
and the volume variation contributes $-\tfrac12(\pi:B)g$.
Together they are \eqref{q4:adj:shift-coefficient}.

The scalar and quadratic terms in \eqref{q4:adj:metric-functional} have
coefficient
\[
 \nabla^2u-(\Delta u)g-u\Ric_g
       +2u(K^2-\tau K)+u\mu g.
\]
Adding \eqref{q4:adj:shift-coefficient} and the test-vector correction therefore
gives
\begin{align}
 0={}&\nabla^2u-(\Delta u)g-u\Ric_g+2u(K^2-\tau K)+u\mu g
       \notag\\
    &+\mathcal L_XK-(\operatorname{div}X)K
      -\{X(\tau)+K^{ab}\nabla_aX_b\}g-X_{(i}J_{j)}.
 \label{q4:adj:metric-adjoint-expanded}
\end{align}
Equation~\eqref{q4:adj:kid-first} gives
$K^{ab}\nabla_aX_b=-u|K|^2$.  Moreover, using
$\nabla^jK_{ij}=J_i+\nabla_i\tau$ in
\eqref{q4:adj:conformal-adjoint} gives
\begin{equation}
 \Delta u+X(\tau)
   =u|K|^2-J(X)=u(|K|^2+\mu).
 \label{q4:adj:trace-identity}
\end{equation}
Together with $\operatorname{div}X=-u\tau$, these identities cancel the scalar
multiples of $g$ in \eqref{q4:adj:metric-adjoint-expanded} and give
\eqref{q4:adj:kid-second}.
\end{proof}

The multiplier now has smooth interior densities and satisfies the
full adjoint system, including its matter term. To use it as a geometric
field we still need its boundary behavior and its normalization at
infinity. Both follow from the same equations and the allowed variations;
neither is imposed as an additional boundary condition.

\subsection{Boundary regularity and the asymptotic constants}

\begin{lemma}[Prolongation and absence of boundary atoms]
\label{q4:adj:boundary-regularity}
The fields $u,X$ extend smoothly to the one-sided boundary $S$.  Their first
jet at one interior point determines them on the connected interior.
The original moment measures have no boundary-supported part, and hence are
exactly $u\,\mathrm dV_g$ and $X\,\mathrm dV_g$ on all of $\Omega$.
\end{lemma}

\begin{proof}
Put $B_{ij}=-uK_{ij}$.  Differentiate
$\nabla_{(i}X_{j)}=B_{ij}$ in three arrangements, add two of the resulting
equations, and subtract the third.  Commuting derivatives gives
\begin{align}
 \nabla_i\nabla_jX_k
 &=\nabla_iB_{jk}+\nabla_jB_{ik}-\nabla_kB_{ij}
       +\mathcal R_{ijk}(X),\label{q4:adj:prolongation-x}\\
 \mathcal R_{ijk}(X)
 &=-\tfrac12\bigl(
       ([\nabla_j,\nabla_i]X)_k
       +([\nabla_i,\nabla_k]X)_j
       +([\nabla_j,\nabla_k]X)_i\bigr).\notag
\end{align}
Every commutator is a curvature contraction with $X$.  Thus this equation
uses only $K\nabla u$, $u\nabla K$, and curvature times $X$.
Expanding the Lie derivative in \eqref{q4:adj:kid-second} gives
\begin{align}
 \nabla_i\nabla_ju
 ={}&u(\Ric_{g,ij}+\tau K_{ij}-2(K^2)_{ij})-X^a\nabla_aK_{ij}
       \notag\\
 &-K_{aj}\nabla_iX^a-K_{ia}\nabla_jX^a+X_{(i}J_{j)}.
 \label{q4:adj:prolongation-u}
\end{align}
Consequently
\[
 \mathcal J=(u,X,\nabla u,\nabla X)
\]
obeys a homogeneous linear first-order system, with smooth coefficients
built from $g,K$, their required derivatives, and $J$.  Along any smooth
curve, this is a linear ODE for $\mathcal J$.  Uniqueness along curves proves
the first-jet assertion.

In a compact collar of $S$, start on an interior collar section and solve
this ODE along the normal segments down to $S$.  Its coefficients are smooth
up to the one-sided boundary.  Smooth dependence on the initial point and on
the normal parameter provides the smooth extension, agreeing with the
original solution for positive collar parameter.  The pointwise inequalities
and complementarity extend by continuity.

The original moments could still have parts supported on $S$; the interior
regularity alone does not exclude them.  Let $s$ be geodesic collar distance,
let $f\in C^\infty(S)$ be arbitrary and extended constantly along the normal
segments, and take a collar cutoff $\chi$ equal to one near $S$.  Choose the
compact pure metric variation
\begin{equation}
 h=\frac{s^2\chi(s)f}{6}\bigl(g-\mathrm ds\otimes\mathrm ds\bigr),
 \qquad p=0.
 \label{q4:adj:boundary-atom-test}
\end{equation}
This is one of the allowed smooth compact variations, with strict-direction
coefficient $a=0$.  Its value and first jet vanish on $S$.  The trace along
the three-dimensional collar slices is $s^2\chi f/2$, so its second normal
derivative at $S$ is $f$.  At $S$, the scalar variation
\[
 R'_g(h)=\nabla^i\nabla^jh_{ij}-\Delta(\tr_gh)-\Ric_g:h
\]
therefore equals $-f$: the only nonzero second jets are normal-normal
derivatives of tangential components, whereas $h_{ss}=h_{sA}=0$.
The variation of $\tau^2-|K|^2$ uses only $h$, that of $J$ at fixed covariant
$K$ uses only $h,\nabla h$, and the test-vector correction uses only $h$.
Consequently $C'_H|_S=-f$, independently of the active ray $v$.

The area derivative vanishes because $h|_S=0$; the expansion derivative
vanishes because $h|_S=\nabla h|_S=p|_S=0$; and the charge derivative
vanishes by compact support.  Integrating the smooth interior densities by
parts gives zero as well: the adjoint vanishes, and its scalar boundary
terms involve only $h,\nabla h$, while its momentum boundary terms involve
only $h$.  This is ordinary integration up to $S$ using the one-sided smooth
fields, not distributional differentiation of a zero extension across $S$.
Thus \eqref{q4:adj:localized-multiplier} says that the scalar boundary measure
annihilates every $f\in C^\infty(S)$, and that measure is zero.
The inequality $|X|_g\le u$ as measures eliminates the vector boundary part.
Moreover, since $u$ is the pushforward of the positive measure $\Lambda$,
the entire measure $\Lambda$ has zero mass over $S$, not merely zero first
moment there.
\end{proof}

We first obtain constant limits from a Hessian estimate and causality.
Related radial estimates for asymptotic Killing fields appear in
\cite[Proposition~2.1 and Appendix~C]{BeigChrusciel1996}.
The normalization of the limits will then follow from the end variations.

\begin{lemma}[Asymptotics of causal fields]
\label{q4:adj:causal-jet-asymptotics}
Let $g$ be a smooth metric uniformly comparable to the Euclidean metric
on an end $\R^4\setminus\overline B_R$. Let $u$ be a smooth function
and $X$ a smooth vector field there with $u\ge |X|_g$. Write $Y$ for
all coordinate components of $(u,X)$. Suppose that, for some $q_0>1$
and a constant $C$, the coordinate derivatives satisfy
\begin{equation}
 |\partial^2Y|
 \le C r^{-1-q_0}|\partial Y|+C r^{-2-q_0}|Y|.
 \label{q4:adj:jet-estimate}
\end{equation}
Then there are constants $u_\infty$ and $X_\infty\in\R^4$ such that
\[
 (u,X)=(u_\infty,X_\infty)+O_2(r^{-q_0}).
\]
\end{lemma}

\begin{proof}
Fix a sufficiently large coordinate sphere of radius $R$ and put
\[
 D(r)=\sup_{\omega\in S^3}|\partial Y(r\omega)|,
 \qquad \mathcal M(r)=\sup_{R\le t\le r}D(t).
\]
Along each ray,
$|Y(t\omega)|\le C+\int_R^tD(a)\,\mathrm da\le C+t\mathcal M(t)$,
uniformly in $\omega$.  Integrating
$\partial_r(\partial Y)$ using \eqref{q4:adj:jet-estimate}, and then taking
the supremum over directions and radii at most $r$, gives
\[
 \mathcal M(r)\le C+C\int_R^r t^{-1-q_0}\mathcal M(t)\,\mathrm dt.
\]
The kernel is integrable.  Gronwall therefore bounds $\mathcal M$ uniformly,
so $\partial Y=O(1)$ and $Y=O(r)$.
Equation~\eqref{q4:adj:jet-estimate} now gives
$\partial^2Y=O(r^{-1-q_0})$.  Each coordinate derivative has a radial limit
$L(\omega)$, with error $O(r^{-q_0})$.  Any two points on the sphere of
radius $r$ can be joined by a spherical arc of length at most $\pi r$;
integrating the Hessian along this arc gives an $O(r^{-q_0})$ difference
between their gradients.  Letting $r\to\infty$ shows that $L(\omega)$ is
one constant matrix $L$.  Radial integration of $\partial Y-L$ then gives
\[
 \partial Y=L+O(r^{-q_0}),\qquad Y=Lx+O(1),
\]
where boundedness of the second error uses $q_0>1$.
Since $u\ge0$ in every direction, the linear part of $u$ vanishes.
Uniform comparability of $g$ with the Euclidean metric and $|X|_g\le u$
then force the linear part of $X$ to vanish as well.  Thus $Y=O(1)$
and $\partial Y=O(r^{-q_0})$.  Substitution in
\eqref{q4:adj:jet-estimate}, using $q_0>1$, improves the Hessian bound to
$O(r^{-2-q_0})$.  Integrating each gradient component radially toward its
zero limit gives $\partial Y=O(r^{-1-q_0})$.  A further radial integration
gives $Y=B(\omega)+O(r^{-q_0})$; comparison of values along the same spherical
arcs now gives an $O(r^{-q_0})$ difference, so $B$ too is independent of
direction. This proves the asserted constant limit and both differentiated
estimates.
\end{proof}

\begin{lemma}[Asymptotic normalization and positive lapse]
\label{q4:adj:asymptotics}
Equation~\eqref{q4:adj:decay} holds, and $u>0$ on
$\operatorname{int}\Omega$.
\end{lemma}

\begin{proof}
Let $Y$ denote all coordinate components of the pair $(u,X)$.  The prolonged
system and the original $O_2$ metric and $O_1$ tensor decay imply
Equation~\eqref{q4:adj:jet-estimate}.
Indeed curvature, $\nabla K$, and $J$ have order $-2-q_0$, while $K$ and
the Christoffel symbols have order $-1-q_0$.  Equations
\eqref{q4:adj:prolongation-x} and \eqref{q4:adj:prolongation-u} use no higher
background derivatives.  Converting their covariant Hessians to coordinate
Hessians adds terms with coefficients $\Gamma$, $\partial\Gamma$, and
$\Gamma^2$, controlled by the stated $O_2$ metric decay.

Lemma~\ref{q4:adj:causal-jet-asymptotics} now gives constants with
\begin{equation}
 (u,X)=(u_\infty,X_\infty)+O_2(r^{-q_0}).
 \label{q4:adj:unnormalized-decay}
\end{equation}

Use the five end prototypes of Definition~\ref{q4:var:variation-space},
whose charge derivatives are computed in Lemma~\ref{q4:var:end-variations}, in
\eqref{q4:adj:localized-multiplier}.  Their strict-direction coefficient is
$a=0$, so the term $za$ is absent; their boundary and area terms vanish
because the prototypes are supported away from $S$.  Integration by parts
using the interior adjoint equations leaves only the outer boundary flux.
By \eqref{q4:adj:unnormalized-decay}, its limit is
\begin{equation}
 6\omega\bigl(u_\infty E'+X_\infty^iP'_i\bigr).
 \label{q4:adj:charge-flux}
\end{equation}
For clarity, the leading integrand is
\[
 u_\infty(\partial_jh_{ij}-\partial_i h_{jj})
 +2X_\infty^j\bigl(p_{ij}-(\tr_\delta p)\delta_{ij}\bigr).
\]
This is exactly \eqref{q4:adj:charge-flux} with the energy and momentum
normalizations of Equations~\eqref{q4:intro:energy}--\eqref{q4:intro:momentum}.  Every other boundary term
contains a decaying background or adjoint coefficient, or its derivative,
paired with $h=O_2(r^{-2})$ or $p=O_1(r^{-3})$ at the corresponding order.
Multiplication by the sphere area $O(r^3)$ makes each such flux tend to zero.
The volume pairing converges as well: the prototype constraint derivatives
are $O(r^{-4-q})$ and the moment densities are bounded.

The scalar prototype $h=r^{-2}\delta$, $p=0$, has $(E',P')=(1,0)$.
For the tensor prototype indexed by a constant vector $e$, its Euclidean
trace reversal contracts with the sphere normal to $r^{-3}e$, and thus has
$(E',P')=(0,e/3)$.  These independent charges identify
$(u_\infty,X_\infty)=(b^0,b)$ from \eqref{q4:adj:localized-multiplier} and
\eqref{q4:adj:charge-flux}.  This proves \eqref{q4:adj:decay}.

If $u$ vanished at an interior point, its nonnegativity would give
$\mathrm du=0$ there.  Taylor's theorem and $|X|\le u$ would give
$X=0$ and $\nabla X=0$ at the same point.  The whole first jet would vanish,
contradicting Lemma~\ref{q4:adj:boundary-regularity} and $u_\infty=b^0>0$.
\end{proof}

\Needspace{8\baselineskip}
\subsection{Boundary fluxes and the two lapse cases}

\begin{lemma}[Boundary equations]
\label{q4:adj:boundary-data}
The multiplier measure at $S$ is $\mathrm d\eta=2u\,\mathrm dA_g$, and
\eqref{q4:adj:horizon-data} holds.
\end{lemma}

\begin{proof}
Take arbitrary compact pure tensor variations up to $S$.  The actual outward
normal of the exterior domain at its inner boundary is $-\nu$.  Integration
by parts in \eqref{q4:adj:localized-multiplier} therefore gives
\[
 -2\int_S\bigl[p(X,\nu)-(\tr_gp)X_\nu\bigr]\,\mathrm dA_g
      -\int_S(\tr_Sp)\,\mathrm d\eta=0.
\]
Arbitrary mixed components $p_{\nu A}$ imply $X_{TS}=0$, and arbitrary
tangential trace implies $\mathrm d\eta=2X_\nu\,\mathrm dA_g$.

Now take $(h,p)=(2\psi g,\psi K)$ with arbitrary compact support up to $S$.
At a MOTS, its expansion derivative is $3\partial_\nu\psi$, while its area
derivative is $3\int_S\psi\,\mathrm dA_g$.  Integrating its constraint
pairing with \eqref{q4:adj:conformal-adjoint} gives
\begin{align*}
 -3c\int_S\psi\,\mathrm dA_g
 ={}&6\int_S\left[u\partial_\nu\psi
       -\{\partial_\nu u+K(X,\nu)\}\psi\right]\,\mathrm dA_g
       -3\int_S\partial_\nu\psi\,\mathrm d\eta.
\end{align*}
The boundary value and normal derivative of $\psi$ can be prescribed
independently.  Their coefficients give $\mathrm d\eta=2u\,\mathrm dA_g$
and $\partial_\nu u+K(X,\nu)=c/2$.  Together with the tensor-variation
conclusions these are exactly \eqref{q4:adj:horizon-data}.
\end{proof}

\begin{lemma}[Boundary lapse dichotomy]
\label{q4:adj:lapse-dichotomy}
Either $u>0$ at every point of $S$ or $u=0$ identically on $S$.
\end{lemma}

\begin{proof}
Let $L_S$ be the expansion variation at $S$ under normal motion with speed
$s\nu$.  Its principal part is $-\Delta_S$.  More explicitly, in a geodesic
normal extension of $\nu$, if $\mathrm{II}$ is the second fundamental form
of $S$ in $(\Omega,g)$,
\[
 L_Ss=-\Delta_Ss+2K(\nu,\nabla_Ss)
  +\left[-|\mathrm{II}|^2-\Ric_g(\nu,\nu)
          +\tr_{TS}(\nabla_\nu K)\right]s.
\]
Only the smoothness of the lower coefficients and the displayed principal
part will be needed.

Extend $s\nu$ to a smooth compactly supported vector field $Y_1$ up to the
boundary, and use $H=(\mathcal L_{Y_1}g,\mathcal L_{Y_1}K)$ in
\eqref{q4:adj:localized-multiplier}.  These are legitimate one-sided tensor
variations; one may compute their jets using any smooth extension across
$S$.  Their area and expansion derivatives are respectively
$\int_SHs\,\mathrm dA_g$ and $L_Ss$.

We check the active-ray derivative with the specified vector identification.
At a fixed interior point put $\mathcal A(v)=\nabla_vY_1$, so
$h^\sharp=\mathcal A+\mathcal A^*$, where the adjoint uses $g$.
Let $\phi_t$ be the local flow of $Y_1$ and let $v_t$ be the isometrically
identified test vector for $g_t=\phi_t^*g$.  Naturality gives
\[
 C_{\phi_t^*g,\phi_t^*K}(x,v_t)
   =C_{g,K}(\phi_t x,\mathrm d\phi_t(v_t)).
\]
At $t=0$, the covariant velocity of the vector on the right, relative to
parallel transport along $\phi_t x$, is
\[
 \mathcal A(v)+\dot v
 =\mathcal A(v)-\tfrac12(\mathcal A+\mathcal A^*)v
 =\tfrac12(\mathcal A-\mathcal A^*)v.
\]
It is perpendicular to $v$.  The base derivative of $C$, with $v$ parallel
transported, vanishes at an interior active ray because it differentiates
a nonnegative function at a two-sided interior minimum.  The remaining
vector derivative is $2J(\tfrac12(\mathcal A-\mathcal A^*)v)$, also zero:
at an active unit ray $J=-\mu v^\flat$, and at an active ray with $|v|<1$
we have $\mu=J=0$.  Therefore $C'_H=0$ at every interior active ray.

The vector field $Y_1$ need not preserve $S$: its compact Lie derivatives
are permitted tensor variations, independently of feasibility of a flow on
the whole exterior.  No normal derivative at a one-sided boundary minimum
is being set equal to zero.  Such boundary rays contribute no integral
because $\Lambda$ has zero mass over $S$ by
Lemma~\ref{q4:adj:boundary-regularity}.  Thus the entire constraint pairing
vanishes.  The multiplier identity and $\mathrm d\eta=2u\,\mathrm dA_g$
now read
\[
 -c\int_S Hs\,\mathrm dA_g=-2\int_SuL_Ss\,\mathrm dA_g,
\]
and hence give
\begin{equation}
 2L_S^*u=cH.
 \label{q4:adj:boundary-stability-adjoint}
\end{equation}
Here the adjoint is with respect to $\mathrm dA_g$ on the closed manifold
$S$.

Outer area minimization, tested against every nonnegative outward normal
speed, gives $H\ge0$ pointwise.  Hence $L_S^*u\ge0$ and $u\ge0$ on $S$.
Increase the zeroth coefficient of $L_S^*$ by a constant until that
coefficient is nonnegative.  The inequality remains valid because $u\ge0$.
The strong minimum principle for this operator with principal part
$-\Delta_S$ shows that a zero of $u$ forces $u$ to vanish on the connected
surface $S$.  Otherwise it is everywhere positive.
\end{proof}

We have identified the field at the horizon and at infinity. These
identities let us construct a stationary spacetime containing the original
slice. Causality and complementarity will make that spacetime vacuum on
the timelike region; an initially possible null stress must remain in the
calculation until that region is identified.

\subsection{The constructed spacetime and its null stress}

\begin{lemma}[Killing development and timelike collars]
\label{q4:adj:killing-development}
The metric in \eqref{q4:adj:stationary-metric} has all the properties in
parts \textup{(iv)} and \textup{(v)} of Theorem~\ref{q4:adj:causal-field}.
\end{lemma}

\begin{proof}
The interior positivity of $u$ makes \eqref{q4:adj:stationary-metric} a smooth
Lorentzian metric, with future normal $n=u^{-1}(\partial_t-X)$ after choosing
the indicated time orientation.  In the convention of the problem, the
second fundamental form of a lapse-shift metric is
\[
 \frac1{2u}\bigl(\partial_tg-\mathcal L_Xg\bigr).
\]
Its coefficients here are independent of $t$, and
\eqref{q4:adj:kid-first} identifies this tensor with the original $K$.
The vector field $\xi=\partial_t$ is Killing by construction.

We record the curvature computation to fix both the matter term and the
sign convention.  The Gauss and normal-variation identities for a general
lapse-shift metric, with positive second fundamental form convention, are
\begin{align*}
 \Ric_{\mathbf g,ij}
  &=\Ric_{g,ij}+\tau K_{ij}-2(K^2)_{ij}
    +u^{-1}\bigl(\partial_tK_{ij}
           -(\mathcal L_XK)_{ij}-(\nabla^2u)_{ij}\bigr),\\
 \Ric_{\mathbf g}(n,n)
  &=-u^{-1}(\partial_t\tau-X\tau)-|K|^2+u^{-1}\Delta u.
\end{align*}
Tracing these equations uses
\[
 g^{ij}\bigl(\partial_tK_{ij}-(\mathcal L_XK)_{ij}\bigr)
 =\partial_t\tau-X\tau+2u|K|^2.
\]
In the stationary case this yields
\begin{align}
 \Ric_{\mathbf g,ij}
 &=\Ric_{g,ij}+\tau K_{ij}-2(K^2)_{ij}
      -u^{-1}\bigl((\mathcal L_XK)_{ij}+(\nabla^2u)_{ij}\bigr),
 \label{q4:adj:spatial-ricci}\\
 R_{\mathbf g}
 &=R_g+\tau^2+|K|^2-2u^{-1}(\Delta u+X\tau).
 \label{q4:adj:spacetime-scalar}
\end{align}
Equation~\eqref{q4:adj:trace-identity} and complementarity show that
\[
 R_{\mathbf g}=2\mu+\frac{2J(X)}u=0.
\]
Writing $\mathbf G$ for the Einstein tensor, the constraints and
\eqref{q4:adj:kid-second} consequently give
\begin{equation}
 \mathbf G(n,n)=\mu,\qquad
 \mathbf G(n,e_i)=J_i,\qquad
 u\mathbf G_{ij}=-X_{(i}J_{j)}.
 \label{q4:adj:einstein-components}
\end{equation}
These identities can also be compared with the transversal Killing
construction in \cite{BeigChrusciel1997KID}; its momentum symbol has the
opposite sign to the one used here.  The null-fluid structure of a modified
constraint adjoint is discussed in \cite[Section~6.1]{HuangLee2024Bartnik}.
We use the displayed direct calculation, in four spatial dimensions, rather
than importing a stationarity or dimensional-transfer theorem.

If $\mu=0$, the dominant energy condition gives $J=0$, so all entries in
\eqref{q4:adj:einstein-components} vanish.  If $\mu>0$, the inequalities
\[
 0=u\mu+J(X)\ge u\mu-|J||X|\ge0
\]
force $|J|=\mu$, $|X|=u$, and $J=-\mu X^\flat/u$.
Since $\mathbf g(\xi,n)=-u$ and $\mathbf g(\xi,e_i)=X_i$, the tensor with
components \eqref{q4:adj:einstein-components} is exactly
\begin{equation}
 \mathbf G=\frac\mu{u^2}\,\xi^\flat\otimes\xi^\flat
 \quad\hbox{where }\mu>0.
 \label{q4:adj:null-stress}
\end{equation}
The Killing one-form in this equation is the spacetime one-form.
It is null wherever the displayed tensor is nonzero.  Thus
$\mathbf G(\xi,\cdot)=0$ everywhere.  Since $R_{\mathbf g}=0$, this is the
first assertion of \eqref{q4:adj:ricci}.  If $N>0$, then $|X|<u$, so
complementarity forces $\mu=J=0$; all the Ricci components vanish.  This proves
the second assertion.  An internal null region has not been excluded or
assumed vacuum.

It remains to examine the original boundary and the end.  Equation~\eqref{q4:adj:decay}
gives $N\to(b^0)^2-|b|^2=1$.  At $S$, Equation~\eqref{q4:adj:horizon-data}
gives $X=u\nu$, so $N=0$ and all tangential derivatives of $N$ vanish.
If $u>0$ on $S$, Equation~\eqref{q4:adj:kid-first} gives
$\langle\nabla_\nu X,\nu\rangle=-uK(\nu,\nu)$.  Hence
\[
 \partial_\nu N
 =2u\bigl(\partial_\nu u+uK(\nu,\nu)\bigr)=2u\kappa,
\]
which is \eqref{q4:adj:positive-collar}.  Compactness of $S$ supplies a full
collar on which $N>0$.

In the other case, $u|_S=0$ and $X|_S=0$.  Tangential derivatives of $X$
vanish.  The normal-normal and normal-tangential components of
\eqref{q4:adj:kid-first} then give $\nabla_\nu X=0$ on $S$, so every first
derivative of $X$ vanishes there.  Smooth one-sided division by normal
distance gives $u=sd$ and $X=s^2Y_2$.  The boundary condition gives
$d|_S=\partial_\nu u|_S=\kappa>0$.  This is
\eqref{q4:adj:zero-collar}, and once again compactness gives a full collar
with $N>0$.  Together with positivity far out on the unique end, these
collars place all interior zeros of $N$ in an interior compact set.
\end{proof}

\begin{proof}[Proof of Theorem~\ref{q4:adj:causal-field}]
Combine Lemmas~\ref{q4:adj:interior-equations},
\ref{q4:adj:boundary-regularity}, \ref{q4:adj:asymptotics},
\ref{q4:adj:boundary-data}, \ref{q4:adj:lapse-dichotomy}, and
\ref{q4:adj:killing-development}.
\end{proof}
